\documentclass[twocolumn, 12pt]{article}

\usepackage[utf8]{inputenc}
\usepackage[english, spanish]{babel}
\usepackage{fullpage}
\usepackage{graphicx}
\usepackage{amsmath}
\usepackage{enumitem}
\usepackage{chngcntr}
\usepackage{setspace}
\usepackage{url}
\usepackage{csquotes}
\usepackage{float}
\usepackage{verbatim}
\usepackage{tabularx}
\usepackage{amsmath}
\usepackage{caption}

\counterwithin{figure}{section}
\renewcommand{\thesection}{\arabic{section}}
\renewcommand{\thesubsection}{\thesection.\arabic{subsection}}
\renewcommand{\baselinestretch}{1.5}

\usepackage[style=apa, maxnames=6, minnames=3, backend=biber]{biblatex}
\DefineBibliographyStrings{english}{%chktex-file 1 chktex-file 6
    andothers = {\em et\addabbrvspace al\adddot}
}
\addbibresource{./Bibliography/bibliography.bib}

\usepackage{array}
\usepackage{enumitem}

\setlength{\parskip}{0pt}

\raggedbottom{}

\begin{document}

\begin{titlepage}
    \centering
    \includegraphics[width=0.3\textwidth]{Images/logo_utb.png}\par\vspace{1cm}
    {\scshape\LARGE Universidad Tecnológica de Bolívar \par}
    \vspace{1cm}

    {\scshape\Large FÍSICA ELÉCTRICA \par}
    \vspace{.2cm}

    % chktex-file 8
    {\scshape\Large M2 - C \par}
    \vspace{1cm}
    % chktex-file 8
    \slshape {\Large \bfseries{}Informe de Laboratorio No. VII\\}
    \vspace{1cm}

    \slshape {\itshape{} Mauro González, T00067622 \\}
    \slshape {\itshape{} German De Armas Castaño, T00068765 \\}
    \slshape {\itshape{} Angel Vega Rodriguez, T00068186 \\}
    \slshape {\itshape{} Juan Jose Osorio Ariza, T00067316 \\}
    \slshape {\itshape{} Juan Eduardo barón, T00065901 \\}
    \vfill
    Revisado Por \\
    Gabriel Hoyos Gomez Casseres\\
    {\large \today\par}
\end{titlepage}

% ----------------------------------------------------------------------|>
\section{Introducción}

% ----------------------------------------------------------------------|>
\section{Objetivos}

\subsection{Objetivo general}

\subsection{Objetivos específicos}

% ----------------------------------------------------------------------|>
\section{Marco Teórico}

% -------------------------------------------------------|>
\subsection{Leyes de Kirchhoff} \nocite{KhanAcademy__Kirchhoff}

% --------------------------------|>
\subsubsection{Ley de Corriente}

La ley de corriente de Kirchhoff establece que la suma
algebraica de todas las corrientes que convergen en un nodo
es igual a la suma algebraica de las corrientes que
divergen del nodo. Matemáticamente, se puede expresar de la
siguiente manera:

{\large
\begin{equation}
    \sum i_{adentro} = \sum i_{afuera}
\end{equation}
}

% --------------------------------|>
\subsubsection{Ley de Voltaje}

La suma de los voltajes alrededor de una malla es igual a
cero, lo cual se puede expresar mediante la ley de voltaje
de Kirchhoff de la siguiente manera:

En su forma general, considerando \textit{n} como el número
de voltajes de los componentes en la malla, se tiene:

{\large
\begin{equation}
    \sum v_n = 0
\end{equation}
}

Asimismo, la ley de voltaje de Kirchhoff se puede enunciar
de la siguiente manera: alrededor de una malla, la suma de
las subidas de voltaje es igual a la suma de las bajadas de
voltaje. Esto se expresa como:

{\large
\begin{equation}
    \sum v_{\text{subida}} = \sum v_{\text{bajada}}
\end{equation}
}

% ----------------------------------------------------------------------|>
\section{Montaje Experimental}

% ----------------------------------------------------------------------|>
\section{Datos Experimentales}

\begin{table}[H]
    \captionsetup{justification=centering}
    \centering

    % chktex-file 44
    \begin{tabularx}{0.9\linewidth}{|>{\centering\arraybackslash}X|>{\centering\arraybackslash}X|>{\centering\arraybackslash}X|>{\centering\arraybackslash}X|}
        \multicolumn{4}{c}{Valor de resistencias $(K\Omega)$} \\ \hline

        $R_1$ & $R_2$  & $R_3$  & $R_4$                       \\ \hline
        $10$  & $0.47$ & $0.33$ & $2.2$                       \\ \hline
    \end{tabularx}

    \caption{Resistencia}

    % chktex-file 24
    \label{tab:datosExperimentales__Resistencias}
\end{table}

\vspace{.5cm}

\begin{table}[H]
    \captionsetup{justification=centering}
    \centering

    \begin{tabularx}{0.9\linewidth}{|>{\centering\arraybackslash}X|>{\centering\arraybackslash}X|>{\centering\arraybackslash}X|>{\centering\arraybackslash}X|>{\centering\arraybackslash}X|}
        \multicolumn{5}{c}{Valor de corrientes $(mA)$} \\ \hline

        $I_1$ & $I_2$ & $I_3$ & $I_4$ & $I_5$          \\ \hline
        $9.9$ & $3.6$ & $6.1$ & $0.5$ & $5.3$          \\ \hline
    \end{tabularx}

    \caption{Corriente}

    % chktex-file 24
    \label{tab:datosExperimentales__Corriente}
\end{table}

\vspace{.5cm}

\begin{table}[H]
    \captionsetup{justification=centering}
    \centering

    \begin{tabularx}{0.9\linewidth}{|>{\centering\arraybackslash}X|>{\centering\arraybackslash}X|>{\centering\arraybackslash}X|>{\centering\arraybackslash}X|}

        \multicolumn{4}{c}{Diferencias de potencial $(V)$}        \\ \hline
        $E1 \linebreak = B_{ab}$ & $V_{bc}$ & $V_{cd}$ & $V_{ef}$ \\ \hline
        $11.92$                  & $-10.06$ & $-1.86$  & $-1.8$   \\ \hline

        $E2 \linebreak = V_{gh}$ & $V_{hf}$ & $V_{ce}$ & \dots    \\ \hline
        $-0.45$                  & $-1.4$   & $0$      & \dots    \\ \hline
    \end{tabularx}

    \caption{Voltaje}

    % chktex-file 24
    \label{tab:datosExperimentales__DIFFPotencial}
\end{table}

\vspace{.5cm}

\begin{table}[H]
    \captionsetup{justification=centering}
    \centering

    \begin{tabularx}{0.9\linewidth}{|>{\centering\arraybackslash}X|>{\centering\arraybackslash}X|>{\centering\arraybackslash}X|>{\centering\arraybackslash}X|>{\centering\arraybackslash}X|}
        \multicolumn{5}{c}{Valor de corrientes $(mA)$} \\ \hline

        $I_1$  & $I_2$  & $I_3$ & $I_4$ & $I_5$        \\ \hline
        $9.98$ & $3.49$ & $5.9$ & $0.5$ & $5.6$        \\ \hline

    \end{tabularx}
    \caption{Corriente \textit{(FEM)}Invertida}

    % chktex-file 24
    \label{tab:datosExperimentales__Corriente-FEMInvertida}
\end{table}

% ----------------------------------------------------------------------|>
\section{Análisis de datos}

% -----------------------------------------------------|>
\subsection{Sume las diferencias de potencial en cada uno de los elementos del circuito para cada
    malla. Registre sus cálculos en la tabla 5.}

\begin{table}[H]
    \captionsetup{justification=centering}
    \centering

    \begin{tabularx}{\linewidth}{|>{\centering\arraybackslash}X|>{\centering\arraybackslash}X|}
        \hline

        \textbf{Malla} & \textbf{V} \\\hline

        $M_1$          &
        {\scriptsize $V_{ab} + V_{bc} + V_{cd}$} $= 11.92 - 10.05 - 1.86 = 0.01$ \linebreak
        $11.92 \simeq 11.91$        \\ \hline

        $M_2$          &
        {\scriptsize $V_{cd} + V_{ef}$} $= 1.86 - 1.8 = 0.06$ \linebreak
        $1.86 \simeq 1.8$           \\\hline

        $M_3$          &
        {\scriptsize $V_{gh} + V_{hf} + V_{ef}$} $= 0.45 - 1.4 + 1.8 = -0.05$ \linebreak
        $1.8 \simeq 1.85$           \\ \hline

        $A_{bgh}$      &
        {\scriptsize $V_{ab} + V_{bg} + V_{gh} + V_{ha}$} $= 11.92 - 10.06 + 0.45 - 1.81 = 0.5$ \linebreak
        $11.92 \simeq 11.87$        \\\hline

        $A_{bef}$      &
        {\scriptsize $V_{ab} + V_{be} + V_{ef}$} $= 11.92 - 10.06 - 1.8 = -0.48$ \linebreak
        $11.92 \sim 12.4$           \\\hline

        $A_{ghd}$      &
        {\scriptsize $V_{cg} + V_{gh} + V_{hd} + V_{dc}$} $= 0 + 0.45 - 1.8 + 1.86 = 0.51$ \linebreak
        $1.8 \sim 2.31$             \\\hline
    \end{tabularx}

    \caption{Diferencias de potencial}

    \label{tab:analisisDatos__1}
\end{table}

% -----------------------------------------------------|>
\subsection{¿Se cumple la ley de las mallas? ¿por qué?}

Si se verifica la ley de malla en un circuito, se observa
que los valores prácticos y teóricos concuerdan, aunque
pueda existir un pequeño error debido a la tolerancia
inherente de los elementos utilizados o a posibles
imprecisiones en los dispositivos de medición.

Es importante destacar que en ciertas ocasiones, este error
puede ser más evidente en ciertas partes del circuito, lo
cual puede requerir una consideración adicional para
evaluar con precisión el grado de coincidencia entre los
valores prácticos y teóricos.

% -----------------------------------------------------|>
\subsection{¿Si no resulta lo que se espera, a qué se deberá?}

Si los resultados obtenidos no son consistentes con las
expectativas, esto puede ser atribuido a diversos factores
de error, tales como mediciones inadecuadas, fallas en los
dispositivos utilizados o errores en los cálculos
realizados.

% -----------------------------------------------------|>
\subsection{¿Si realiza el recorrido en sentido contrario, también se cumple la ley de las mallas?
    ¿Cuál es la diferencia?}

La ley de la corriente de Kirchhoff establece que si se
cumple la conservación de carga en un nodo, la suma de las
corrientes que entran al nodo es igual a la suma de las
corrientes que salen del nodo. En caso de que no se cumpla
esta condición, las corrientes resultantes tendrán la misma
magnitud pero con signos opuestos para indicar la dirección
opuesta del flujo de corriente.

% -----------------------------------------------------|>
\subsection{¿Por qué $V_{ce}$ es cero?}

$V_{ce}$ es cero en un cable debido a que los extremos del cable
están directamente conectados, lo que implica que no hay
una diferencia de potencial significativa entre ellos. En
otras palabras, el voltaje en ambos extremos del cable es
igual pero con signos opuestos, lo cual resulta en un valor
de cero para Vce.

% -----------------------------------------------------|>
\subsection{Sume las corrientes que salen en cada nodo. Registre sus datos en la tabla 6}

Según los datos en la tabla~\ref{tab:datosExperimentales__Corriente},

\begin{table}[H]
    \captionsetup{justification=centering}
    \centering

    \begin{tabularx}{0.9\linewidth}{|>{\centering\arraybackslash}X|>{\centering\arraybackslash}X|>{\centering\arraybackslash}X|}
        \hline
        \textbf{Nodo} & \textbf{Entra}               & \textbf{Sale}                \\\hline

        \textit{C}    & $I_1 = 9.9$                  & $I_2 + I_3 \linebreak = 9.7$ \\\hline

        \textit{E}    & $I_3 = 6.1$                  & $I_4 + I_5 \linebreak = 5.8$ \\\hline

        \textit{D}    & $I_4 + I_5 \linebreak = 5.8$ & $I_3 = 6.1$                  \\\hline

        \textit{F}    & $I_2 + I_3 \linebreak = 9.7$ & $I_1 = 9.9$                  \\\hline

    \end{tabularx}
    \caption{}

    \label{tab:analisisDatos__6}
\end{table}

% ----------------------------------------------------------------------|>
\section{Conclusiones}

\printbibliography

\end{document}